% Corrected French for printed p.189; departures are append-only below.
\clearpage
\SourcePage{189}

\LayerTarget{tohoku:eq:4.2.8}
\LayerTarget{tohoku:proof:theorem-4.2.2}
$A$ étant fixé, pour que ces homomorphismes soient des isomorphismes quels que
soient $p$ et $B$, il faut et il suffit que (i) ce soit le cas pour $p=0$
(ii) $\mathrm{Ext}_{O(x)}^p(A(x),B(x))=0$ si $p>0$ et $B$ est un $O$-module
injectif. Les propositions 4.1.1 et 4.1.2 ont visiblement été faites exprès
pour vérifier ces conditions. On obtient par suite :

\LayerTarget{tohoku:theorem:4.2.2}
\noindent\textsc{Théorème 4.2.2.} \textit{Supposons que $O$ soit un faisceau
cohérent d'anneaux noethériens à gauche, et que $A$ soit un $O$-module
cohérent. Alors les homomorphismes (4.2.8) sont des isomorphismes quels que
soient $B$ et $p$.}

\LayerTarget{tohoku:remark:locally-free-ext-vanishing}
Notons enfin le cas trivial suivant où les homomorphismes (4.2.8) sont des
isomorphismes, les deux membres étant nuls :

\LayerTarget{tohoku:proposition:4.2.3}
\noindent\textsc{Proposition 4.2.3.} \textit{Supposons que $A$ soit localement
isomorphe à $O^n$ (où $n$ est un entier fini $>0$ donné). Alors
$\mathbf{Ext}_O^p(A,B)=0$ pour tout $O$-module $B$ et $p>0$.}

\LayerTarget{tohoku:proof:proposition-4.2.3}
En effet, en vertu de la nature locale de $\mathbf{Ext}_O^p(A,B)$, on peut
supposer $A=O^n$. Mais alors $\mathbf{Hom}_O(A,B)=B^n$ est un foncteur
\textit{exact} en $B$, d'où la conclusion.

\LayerTarget{tohoku:corollary:4.2.3}
\noindent\textsc{Corollaire.} \textit{Si $A$ est localement isomorphe à
$O^n$, on a}
\[
\mathrm{Ext}_O^p(X;A,B)=H^p(X,\mathbf{Hom}_O(A,B))
\]
\textit{pour tout $O$-module $B$.}

\LayerTarget{tohoku:proof:corollary-4.2.3}
Cela résulte en effet aussitôt de la suite spectrale du théorème 4.2.1. En
particulier, on trouve (faisant $A=O$) :

\LayerTarget{tohoku:eq:4.2.9}
\[
\tag{4.2.9}
H^p(X,B)=\mathrm{Ext}_O^p(X;O,B)
\]

\LayerTarget{tohoku:remark:cohomology-as-ext}
formule valable sans aucune hypothèse sur $O$ ou le $O$-module $B$, et qui est
d'ailleurs triviale directement, les $H^p(X,B)$ étant les foncteurs dérivés,
\textit{dans $C^O$} (grâce à prop. 4.1.3) du foncteur
$\Gamma(B)=\mathrm{Hom}_O(O,B)$.

\LayerTarget{tohoku:construction:computing-spectral-sequence}
Pour terminer ce N$^\circ$, nous indiquons une ``méthode de calcul'' de la
suite spectrale du théorème 4.2.1, plus commode que celle qui résulterait de
l'application pure et simple de la définition.

\LayerTarget{tohoku:proposition:4.2.4}
\noindent\textsc{Proposition 4.2.4.} \textit{Soit
$\mathbf{L}=(L_i)$ une résolution gauche de $A$ par des $O$-modules $L_i$
localement isomorphes à des $O^{n_i}$, considérons le foncteur
$\mathbf{h}_{\mathbf{L}}$ sur $C^O$, à valeurs dans les complexes à degrés
positifs dans $C$, défini par
$\mathbf{h}_{\mathbf{L}}(B)=\mathbf{Hom}_O(\mathbf{L},B)$; posons de même
$\mathbf{h}_A(B)=\mathbf{Hom}_O(A,B)$. Alors $\mathbf{h}_{\mathbf{L}}$ est une
résolvante (cf. 2.5) du foncteur $\mathbf{h}_A$.}

\LayerTarget{tohoku:proof:proposition-4.2.4}
En effet, $\mathbf{h}_{\mathbf{L}}$ est un foncteur exact (prop. 4.2.3), de
plus $H^0(\mathbf{h}_{\mathbf{L}}(B))=\mathbf{h}_A(B)$ grâce au fait que le
foncteur $C\longmapsto\mathbf{Hom}_O(C,B)$ est exact à gauche en $C$, enfin si
$B$ est injectif, $\mathbf{h}_{\mathbf{L}}(B)$ est un complexe acyclique en
dimensions $>0$, puisque alors le foncteur
$C\longmapsto\mathbf{Hom}_O(C,B)$ est exact.

\LayerTarget{tohoku:corollary:4.2.4:1}
\noindent\textsc{Corollaire 1.} \textit{Sous les conditions de la proposition
précédente, on a}

\TohokuContinuation{190}
